\documentclass[11pt]{article}

\usepackage[margin=1in]{geometry}
\usepackage{booktabs}
\usepackage{array}
\usepackage{longtable}
\usepackage{listings}
\usepackage{xcolor}
\usepackage{hyperref}
\usepackage{enumitem}
\usepackage{parskip}

\definecolor{codebg}{rgb}{0.96,0.96,0.96}
\lstset{
  basicstyle=\ttfamily\small,
  backgroundcolor=\color{codebg},
  breaklines=true,
  frame=single,
  columns=fullflexible,
  showstringspaces=false
}

\hypersetup{
  colorlinks=true,
  linkcolor=blue,
  urlcolor=blue
}

\title{TCP Window Scaling Attack \\ \large Implementation and Observation Report}
\author{Hillol}
\date{\today}

\begin{document}
\maketitle

\begin{abstract}
This report documents the implementation and live execution of a TCP Window
Scale desynchronization attack against a controlled 3-node lab, together with
its low-bandwidth fallback (Slow Read / zero-window abuse), and three defense
mechanisms tested against the same attack pipeline. For each of the three
attacks demonstrated -- \emph{deflate} (shift $7\to0$ family), \emph{inflate}
(shift $\to14$), and \emph{Slow Read DoS} -- and each of the three defenses
implemented -- static ARP hardening, IPsec AH (transport mode), and a live
WScale-mismatch monitor -- this report states the process actually followed,
the outcome that was expected going in, the outcome that was actually
observed (or, for defenses awaiting a final confirmation run, the outcome
predicted from the mechanism together with an explicit note that it is
pending live verification), an explanation wherever expectation and
observation diverged, and the assumptions the attack and the lab depend on.
\end{abstract}

\tableofcontents

% ---------------------------------------------------------------------------
\section{Introduction}

TCP's Window Scale option (RFC~7323) multiplies the 16-bit advertised receive
window by $2^{\text{shift}}$, where shift $\in [0,14]$. Critically, the shift
value is negotiated \emph{only} once, inside the SYN and SYN-ACK of the initial
three-way handshake, and is never re-advertised or re-verified for the
lifetime of the connection. An on-path attacker who can rewrite that single
option byte in the SYN -- without touching anything else in the packet -- can
make the two endpoints permanently disagree about how large the peer's
receive window really is. This report covers the practical implementation of
that idea and two attack variants built on it, executed against a live lab
built for this project, along with a low-bandwidth fallback attack (Slow
Read) that achieves a similar denial-of-service outcome through a different
mechanism (a deliberately tiny \emph{real} window rather than a
\emph{misread} one).

Three attacks are covered:
\begin{enumerate}
  \item \textbf{Deflate}: the attacker rewrites the client's advertised shift
        \emph{downward} (observed as $10\to0$), so the server believes the
        client's window is far smaller than it really is.
  \item \textbf{Inflate}: the attacker rewrites the client's advertised shift
        \emph{upward} (observed as $10\to14$, and later $1\to14$), so the
        server believes the client's window is far larger than it really is.
  \item \textbf{Slow Read / zero-window DoS} (fallback): no packet tampering
        at all -- many connections are opened and read a few bytes at a time
        with a tiny real receive buffer, tying up the server's connection
        pool.
\end{enumerate}

% ---------------------------------------------------------------------------
\section{Lab Environment}
\label{sec:lab}

All three attacks were executed against a self-contained lab built from
Linux network namespaces (standing in for the deck's ``3 VMs'') connected by
a single Linux bridge acting as the shared switch:

\begin{lstlisting}
[client netns 10.0.0.1] --veth--\
                                  wscale-br0 (bridge/"Switch") --veth-- [attacker netns 10.0.0.3]
[server netns 10.0.0.2] --veth--/
\end{lstlisting}

\begin{itemize}
  \item \textbf{Client} (\texttt{10.0.0.1}) and \textbf{Server}
        (\texttt{10.0.0.2}) are ordinary TCP endpoints with no awareness of
        the attacker.
  \item \textbf{Attacker} (\texttt{10.0.0.3}) is a single-homed Linux host on
        the same broadcast segment, running: \texttt{arp\_spoof.py} (ARP
        poisoning), an \texttt{iptables} \texttt{NFQUEUE} divert rule, and
        \texttt{wscale\_attack.py} (the packet rewrite, via
        \texttt{scapy} + \texttt{NetfilterQueue}).
  \item Throughput and TCP-level impact were measured directly from the
        kernel's \texttt{TCP\_INFO} socket option (retransmits, congestion
        window, RTT, negotiated window scale) rather than relying on
        \texttt{iperf3}, giving byte-for-byte agreement between what each
        endpoint believed and what was measured.
  \item All traffic between client and server was also captured on the
        bridge with \texttt{tcpdump} and inspected both with a custom
        detector script (\texttt{analysis/detect\_wscale\_mismatch.py}) and
        manually in Wireshark.
\end{itemize}

% ---------------------------------------------------------------------------
\section{Attack 1: Deflate (Window Scale Downgrade)}
\label{sec:deflate}

\subsection{Attack Process}
\begin{enumerate}
  \item \textbf{Baseline capture.} A clean 20\,MB transfer between client and
        server was captured and measured with no attacker on-path, to have a
        reference point.
  \item \textbf{Gain MITM.} \texttt{arp\_spoof.py} sent periodic forged ARP
        replies telling the client ``the server is at the attacker's MAC''
        and the server ``the client is at the attacker's MAC.'' With IP
        forwarding enabled on the attacker and no route back out a second
        interface (single NIC), the kernel naturally re-ARPs and relays
        traffic back out the same interface to the real peer once it decides
        to forward it -- the classic single-host ARP-spoofing MITM.
  \item \textbf{Divert.} An \texttt{iptables} rule in the attacker's mangle
        \texttt{FORWARD} chain sent every SYN/SYN-ACK (matched on the SYN
        flag bit alone, so both directions qualify) to \texttt{NFQUEUE}
        queue 0:
\begin{lstlisting}
iptables -t mangle -A FORWARD -p tcp --tcp-flags SYN SYN \
    -j NFQUEUE --queue-num 0 --queue-bypass
\end{lstlisting}
  \item \textbf{Rewrite.} \texttt{wscale\_attack.py}, bound to that queue,
        inspected each queued packet, and -- restricting itself to the
        client's \emph{plain} SYN only (\texttt{--target client-syn}, matching
        the deck's timing diagram) -- located the \texttt{WScale} TCP option
        and overwrote its single shift byte from the client's real value to
        \texttt{0}. Both the TCP and IP checksums were deleted so
        \texttt{scapy} recomputed them on serialization; nothing else in the
        packet (length, sequence number, addresses) was touched, so no
        re-segmentation was necessary.
  \item \textbf{Measure.} A 20\,MB bulk transfer was run again, this time
        with the MITM and rewrite active, and its \texttt{TCP\_INFO} and
        wall-clock throughput were compared against the baseline.
  \item \textbf{Verify.} The bridge-side pcap was run through
        \texttt{detect\_wscale\_mismatch.py} and inspected in Wireshark for
        the tell-tale duplicate SYN.
\end{enumerate}

\subsection{Expected Outcome}
Per the design report: the server ends up believing the client's window is
$2^{\text{real shift}}$ times smaller than it actually is. Since the
client's real shift was 10, a rewrite to shift 0 means the server
under-estimates the client's window by a factor of $2^{10}=1024$. The server
should then throttle its sending to what it believes is the client's tiny
real capacity, and throughput should collapse toward the pre-scaling 64\,KB
regime -- visible as a large, sustained drop in throughput with no actual
packet loss (this is a \emph{self-throttling}, not a \emph{loss-driven}
degradation).

\subsection{Actual Outcome}
The rewrite was confirmed in the attacker's own log:
\begin{lstlisting}
2026-09-19 00:20:48,239 listening on NFQUEUE 0 (target=client-syn, new_shift=0)
2026-09-19 00:20:50,989 rewrote WScale 10.0.0.1:44992 -> 10.0.0.2:9000  shift 10 -> 0
\end{lstlisting}

\begin{table}[h]
\centering
\caption{Deflate attack: baseline vs.\ attack, measured via \texttt{TCP\_INFO}}
\begin{tabular}{lrrrr}
\toprule
Run & Throughput (MB/s) & \texttt{total\_retrans} & RTT ($\mu$s) & Server's belief (\texttt{snd\_wscale}) \\
\midrule
Baseline & 1781.163 & 0 & 360 & 10 (correct) \\
Attack (deflate) & 2.314 & 0 & 9 & 0 (wrong; real is 10) \\
\bottomrule
\end{tabular}
\end{table}

\noindent This is a \textbf{99.9\% throughput drop}, with \emph{zero}
retransmissions -- confirming the collapse is pure self-throttling by the
server based on a wrong belief, not packet loss. The bridge-side capture
showed the predicted duplicate SYN:
\begin{lstlisting}
10.0.0.1:44992 -> 10.0.0.2:9000  seq=2098777608  [SYN]
    wscale values seen: [0, 10]  <-- MISMATCH (tampered in flight)
10.0.0.2:9000 -> 10.0.0.1:44992  seq=389760085  [SYN-ACK]
    wscale values seen: [10]

996 anomalously small effective windows (< 16 KiB) on data packets,
almost all at exactly 63 bytes effective window.
\end{lstlisting}
In Wireshark, the two SYN frames with identical IP addresses and sequence
number but different source MAC addresses (client's real MAC for the first,
attacker's MAC for the re-injected second) are visible a frame or two apart,
each with a different value in \emph{TCP $\to$ Options $\to$ Window Scale}.

\subsection{Discrepancies and Explanation}
The core prediction (large throughput collapse, mismatch visible, small
effective windows) matched expectation exactly and needed no correction.
Two things were observed that the design report did not explicitly call
out, and are worth recording:

\begin{enumerate}
  \item \textbf{Packet volume.} An early, uncapped capture of this same
        scenario recorded \textbf{over 1.3 million packets} for a single
        20\,MB transfer -- far more than a normal transfer would ever
        produce. This is fully explained, not a bug: with an effective
        window collapsed to $\sim\!63$ bytes, moving 20\,MB requires
        roughly $20{,}971{,}520 / 63 \approx 333{,}000$ data segments, each
        needing its own ACK before the next can be sent, and \emph{each of
        those frames crosses the bridge twice} because of the two-hop MITM
        path (client$\to$attacker, then attacker$\to$server, as two distinct
        Ethernet frames sharing the same IP 5-tuple). $333{,}000 \times 2
        \text{ (data+ack)} \times 2 \text{ (hops)} \approx 1{,}332{,}000$,
        which matches the observed count almost exactly. This is a
        genuine, if unstated, consequence of a successful deflate attack:
        it is not merely slow, it is also extremely ``chatty.''
  \item \textbf{Congestion window behaviour.} \texttt{snd\_cwnd} stayed at
        its initial value (10 segments) for the entire transfer rather than
        growing via slow start. This is consistent with the connection being
        \emph{receive-window-limited} rather than congestion-limited: the
        server's window ceiling is far below what would be needed to
        exercise congestion control at all, so cwnd growth (which requires
        filling the current window) never has an opportunity to trigger.
\end{enumerate}

% ---------------------------------------------------------------------------
\section{Attack 2: Inflate (Window Scale Overstatement)}
\label{sec:inflate}

\subsection{Attack Process}
Identical pipeline to Section~\ref{sec:deflate}, except
\texttt{wscale\_attack.py} was run with \texttt{--new-shift 14} instead of
\texttt{0}, so the client's real shift is rewritten \emph{upward} rather than
downward. In addition, after the first attempt (below) failed to show any
effect, the process was revised to also pin the client's \emph{real} socket
receive buffer to a fixed, small size before connecting
(\texttt{throughput/bulk\_client.py --rcvbuf 65536}), which disables Linux's
automatic receive-buffer growth for that socket.

\subsection{Expected Outcome}
The server should believe the client's window is $2^{14-\text{real shift}}$
times \emph{larger} than it actually is. Data that the server believes fits
inside the client's window can therefore exceed the client's \emph{real}
buffer capacity; bytes that arrive beyond what the client can actually
buffer are dropped at the receiver, forcing the sender into retransmissions
-- i.e., a retransmission storm and/or stalls, rather than the pure collapse
seen in the deflate case.

\subsection{Actual Outcome}
\paragraph{First attempt (no pinned receive buffer).} The rewrite was
confirmed (\texttt{shift 10 -> 14}, server's \texttt{snd\_wscale} became
14), but the measured result \emph{diverged from the prediction}: throughput
\emph{increased} slightly (1176.326 $\to$ 2480.288 MB/s) and
\texttt{total\_retrans} stayed at \textbf{0} -- no visible damage at all.

\paragraph{Second attempt (client receive buffer pinned to 65536 bytes).}
With the real buffer fixed and small, the predicted effect appeared clearly
and grew more pronounced across repeated runs against the same peer:

\begin{table}[h]
\centering
\caption{Inflate attack, with pinned client receive buffer: three successive runs}
\begin{tabular}{lrrrrl}
\toprule
Run & Throughput (MB/s) & \texttt{total\_retrans} & \texttt{snd\_cwnd} & Rewrite & Real shift seen \\
\midrule
1st & 47.354 & 306 & 156 & $10\to14$ & 10 \\
2nd & 22.747 & 1065 & 121 & $10\to14$ & 10 \\
3rd & 13.535 & 1398 & 167 & $1\to14$ & \textbf{1} (see below) \\
\bottomrule
\end{tabular}
\end{table}

In the third run, the pcap-level detector reported an even more extreme
result:
\begin{lstlisting}
10.0.0.1:59258 -> 10.0.0.2:9000  seq=20094852  [SYN]
    wscale values seen: [1, 14]  <-- MISMATCH (tampered in flight)

Anomalously small effective windows (< 16 KiB) on data packets:
  advertised_window=0  wscale=14  effective=0 bytes   (x4)
\end{lstlisting}

\subsection{Discrepancies and Explanation}
Two distinct divergences from the design report's prediction were observed
and are worth documenting in detail:

\begin{enumerate}
  \item \textbf{No effect at all on the first attempt.} The rewrite
        \emph{did} happen (confirmed in the log and in
        \texttt{TCP\_INFO}), but no retransmissions occurred and throughput
        was unaffected (indeed slightly higher). \textbf{Cause:} on this
        lab's fast, effectively zero-RTT loopback-speed link, Linux's
        receive-buffer autotuning grows the client's real socket buffer to
        several megabytes very quickly, and the whole 20\,MB transfer
        completed in well under a second. There was neither enough time nor
        enough bytes-in-flight for the server's inflated belief to ever
        actually exceed the client's (auto-grown, generous) real buffer.
        The attack was correctly executed but the \emph{lab conditions}
        (unrealistically fast, low-latency link; short transfer) hid its
        consequence. This was corrected by explicitly pinning the client's
        real receive buffer to a small, fixed value with \texttt{SO\_RCVBUF},
        which disables the autotuning and gives the server's false belief an
        actual, fixed ceiling to overrun -- after which the predicted
        retransmissions appeared immediately.
  \item \textbf{The real shift value changed between runs, unprompted
        ($10\to1$).} The design report implicitly assumes the client's real
        window-scale shift is fixed and known; in practice, across repeated
        test connections between the \emph{same} two IP addresses, the
        Linux kernel's per-destination TCP metrics cache remembered
        parameters (such as effective window/RTT behaviour) from the
        earlier, artificially degraded connections and used them to pick a
        smaller initial window scale for a later, brand-new connection. This
        is a real Linux TCP stack behaviour, not an attack artifact, but it
        meaningfully changed the attack's effective multiplier from
        $2^{14-10}=16\times$ to $2^{14-1}=8192\times$ between runs, and is
        the direct cause of the escalating \texttt{total\_retrans} (306
        $\to$ 1065 $\to$ 1398) and the eventual observed \emph{zero}
        effective window across three otherwise-identical repetitions of
        the same command.
\end{enumerate}

% ---------------------------------------------------------------------------
\section{Attack 3 (Fallback): Slow Read / Zero-Window DoS}
\label{sec:slowread}

\subsection{Attack Process}
This attack does not tamper with any packet. A target HTTP server
(\texttt{server/http\_server.py}) was started with a bounded worker pool
(\texttt{--max-conns 20}), serving a large resource at \texttt{/bigfile} and
a fast \texttt{/health} endpoint. The attack client
(\texttt{slowread/slow\_read\_client.py}) opened 50 concurrent connections,
each requesting \texttt{/bigfile}, each with its socket's
\texttt{SO\_RCVBUF} pinned to 8 bytes, and each reading only a few bytes
every second for 30 seconds. Because the client's real advertised window is
genuinely tiny, the server's \texttt{send()} calls block waiting for TCP
flow control to allow more data through, tying up one worker thread per
connection for the full 30 seconds. \texttt{/health} was polled before,
during, and after the attack to observe the server's availability to a
legitimate client.

\subsection{Expected Outcome}
With 50 slow connections against a 20-connection pool, the pool should be
exhausted, and legitimate requests (the \texttt{/health} check) should be
denied or stall for as long as the attack connections are held open; the
server should recover immediately once those connections close.

\subsection{Actual Outcome}
The outcome matched the prediction exactly:
\begin{lstlisting}
[*] baseline health check on 10.0.0.2:8000 before the attack...
    -> ok=True rtt=0.001s

[*] health check WHILE the slow-read attack is running...
    -> ok=False rtt=3.000s body='[Errno 104] Connection reset by peer'
       <-- DENIED / degraded, pool likely exhausted

[*] summary: 20/50 connections held open for the full duration.

[*] health check AFTER attacker connections released...
    -> ok=True rtt=0.001s
\end{lstlisting}
Exactly 20 of the 50 attack connections were held open for the full 30
seconds -- matching the pool size precisely -- while the remaining 30 were
rejected outright once the pool filled (visible in the server's own log as
repeated \texttt{pool exhausted (20 active) -- rejecting ...} lines) or
reset shortly after. The health check was denied only during the attack
window and recovered to sub-millisecond response time immediately
afterward.

\subsection{Discrepancies and Explanation}
No discrepancy was observed for this attack; the outcome matched the
expected mechanism exactly on the first run. This is consistent with the
design report's own framing of Slow Read as the ``guaranteed-working
backup'' -- unlike the window-scale attacks, it does not depend on any
kernel-level buffer-autotuning behaviour or on the two endpoints having a
particular history with each other, only on the server having a finite
worker/connection pool, which is a very common and simple property of real
servers.

% ---------------------------------------------------------------------------
\section{Defense Mechanisms}
\label{sec:defense}

Per the design report's ``Defense Ideas'' slide, three defenses were
implemented and wired up to run against the exact same attack pipeline used
in Sections~\ref{sec:deflate}--\ref{sec:slowread}, rather than only
described: stopping the MITM position itself (static ARP entries,
approximating switch-level Dynamic ARP Inspection), enforcing
transport/network-layer integrity (IPsec AH, transport mode), and detection
(a live WScale-mismatch monitor, functionally equivalent to the deck's
suggested Zeek/Snort rule). Each is implemented as an executable script
under \texttt{defense/} together with \texttt{defense/run\_defense\_demo.sh},
which re-runs the attack pipeline with one defense active at a time.

\emph{A note on scope.} Zeek was not available through this system's default
package repositories (it ships its own separate per-distribution repository,
which was not added for this project), and Snort/Suricata's default rule
language has no built-in way to compare a TCP option's raw value across two
different packets of the same flow -- exactly what a WScale-mismatch check
needs -- without custom scripting. Rather than ship an untested rule, the
detection logic already validated in Sections~\ref{sec:deflate}
and~\ref{sec:inflate} (\texttt{analysis/detect\_wscale\_mismatch.py}) was
adapted to run live against a capture interface
(\texttt{defense/live\_ids.py}); a reference Zeek script with equivalent
logic is included at \texttt{defense/wscale\_ids.zeek} for documentation,
though it was not executed in this lab.

\subsection{Defense A: Stopping the MITM (Static ARP Entries)}
\label{sec:defense-arp}

\paragraph{Process.} \texttt{defense/harden\_arp.sh} reads each endpoint's
real MAC address directly from its own interface
(\texttt{/sys/class/net/eth0/address}) and installs it into the \emph{other}
endpoint's kernel neighbor table as a \texttt{PERMANENT} entry -- e.g.\ the
client is told ``\texttt{10.0.0.2} is at $\langle$server's real MAC$\rangle$,
permanently.'' Linux does not update a \texttt{PERMANENT} neighbor-table
entry in response to unsolicited or gratuitous ARP traffic, which is
precisely the mechanism \texttt{attacker/arp\_spoof.py} depends on. With the
defense applied, the deflate attack pipeline from
Section~\ref{sec:deflate} was re-run unmodified via
\texttt{sudo bash defense/run\_defense\_demo.sh arp}.

\paragraph{Expected Outcome.} \texttt{arp\_spoof.py}'s forged replies should
be silently ignored by both endpoints, so the attacker should never actually
receive client$\leftrightarrow$server traffic. \texttt{wscale\_attack.py}
should therefore log \emph{zero} rewrites, and the bulk transfer should
complete at essentially baseline throughput with the negotiated window
scale unchanged, matching Section~\ref{sec:deflate}'s pre-attack baseline
rather than its throttled result.

\paragraph{Actual Outcome.} The rewrite log shows the NFQUEUE handler bound
and listening, but with \textbf{no} \texttt{rewrote WScale} line for the
entire run -- it never received a single packet to tamper with:
\begin{lstlisting}
2026-09-19 02:01:56,938 listening on NFQUEUE 0 (target=client-syn, new_shift=0)
(no further lines -- zero rewrites)
\end{lstlisting}
The bridge-side pcap confirms the connection went directly client
$\leftrightarrow$ server with no duplicate/tampered SYN:
\begin{lstlisting}
10.0.0.1:53470 -> 10.0.0.2:9000  seq=1951765398  [SYN]
    wscale values seen: [10]
10.0.0.2:9000 -> 10.0.0.1:53470  seq=24270297  [SYN-ACK]
    wscale values seen: [10]

0 flow(s) with a WScale mismatch, 0 anomalously small window packet(s).
No tampering signatures found in this capture.
\end{lstlisting}
And the transfer itself completed at full, untampered throughput with the
correct window scale on both sides:
\begin{lstlisting}
server: throughput=2196.388 MB/s  total_retrans=0  snd_wscale=10
client: throughput=1827.000 MB/s  total_retrans=0  snd_wscale=10  rcv_wscale=10
\end{lstlisting}

\paragraph{Does it work as expected?} \textbf{Yes, exactly.} The static
(permanent) ARP entries held throughout \texttt{arp\_spoof.py}'s entire
poisoning loop; the attacker never received a single client/server packet to
divert or tamper with, and the connection behaved identically to the
untampered baseline in Section~\ref{sec:deflate}. This is a complete,
unambiguous block of the attack at its very first step -- the attacker
never even achieves the on-path position the rest of the attack chain
depends on.

\subsection{Defense B: Transport/Network-Layer Integrity (IPsec AH, Transport Mode)}
\label{sec:defense-ipsec}

\paragraph{Process.} \texttt{defense/setup\_ipsec\_ah.sh} configures
matching IPsec Security Associations and Policies directly via Linux's
\texttt{ip xfrm} interface in both the client and server namespaces, using
AH (Authentication Header) in transport mode with a fixed, manually
pre-shared HMAC-SHA1-96 key (no IKE daemon -- this demonstrates the
authentication mechanism itself, not a production key-exchange setup). AH
authenticates the entire IP payload, TCP header included, via an Integrity
Check Value (ICV) computed from the shared key. With this active (and
static-ARP hardening \emph{not} active, so the attacker remains on-path),
the same attack pipeline was re-run via
\texttt{sudo bash defense/run\_defense\_demo.sh ipsec}.

\paragraph{Expected Outcome.} The attacker can still recompute the IP/TCP
checksums after editing the WScale option, exactly as in every attack in
this report, but cannot recompute a valid AH ICV without the shared key. Two
outcomes are plausible, and both would count as the defense working:
\begin{enumerate}[label=(\alph*)]
  \item the server's kernel receives the tampered, AH-protected SYN,
        recomputes the expected ICV, finds it does not match, and silently
        drops the packet before it ever reaches the TCP stack -- the
        connection should stall while the client keeps retransmitting a SYN
        that keeps failing the same way; or
  \item since AH transport mode changes the IP header's protocol field from
        TCP (6) to AH (51), the attacker's own
        \texttt{iptables -p tcp --tcp-flags SYN SYN} diversion rule may
        simply stop matching the now-AH-wrapped packets at all, in which
        case \texttt{wscale\_attack.py} never sees them and logs no
        rewrites -- the attacker loses its diversion mechanism entirely,
        not only its tampering.
\end{enumerate}
Either way, the bulk transfer is expected to fail to complete within the
pipeline's 10-second timeout, in clear contrast to Section~\ref{sec:deflate},
where the identical attack against unprotected traffic completed
successfully (just badly throttled).

\paragraph{Actual Outcome.} The SAs and policies installed correctly on both
sides (confirmed via \texttt{ip xfrm state/policy show}). With the attacker
still on-path (ARP spoofing unaffected by IPsec), the rewrite log again
shows the handler listening but logging \textbf{zero} rewrites for the
entire run:
\begin{lstlisting}
2026-09-19 02:23:40,801 listening on NFQUEUE 0 (target=client-syn, new_shift=0)
(no further lines -- zero rewrites)
\end{lstlisting}
The pcap-based detector found \textbf{no SYN or SYN-ACK packets at all} in
the capture (not merely zero mismatches):
\begin{lstlisting}
=== SYN / SYN-ACK WScale values per flow ===
(empty)
0 flow(s) with a WScale mismatch, 0 anomalously small window packet(s).
\end{lstlisting}
And, unlike the prediction above, the transfer \textbf{completed
successfully} rather than timing out -- with the window scale unmodified on
both sides, but with visible secondary degradation (202 retransmits on the
server side, throughput down to $\sim$156--158\,MB/s from an unprotected
baseline in the thousands of MB/s):
\begin{lstlisting}
server: throughput=158.140 MB/s  total_retrans=202  snd_wscale=10 (unchanged)
client: throughput=156.081 MB/s  total_retrans=0    snd_wscale=10 (unchanged)
\end{lstlisting}

\paragraph{Does it work as expected?} \textbf{Yes, on the point that
matters, but through outcome (b) specifically, and one detail of the
prediction was wrong.} The attack was completely neutralized -- zero
rewrites, zero mismatches, window scale untouched on both ends -- confirming
outcome (b): AH's protocol-51 wrapping made the attacker's
\texttt{--tcp-flags SYN SYN} diversion rule stop matching the SYN entirely,
so \texttt{wscale\_attack.py} never even saw a packet to tamper with.
Outcome (a) (an AH ICV rejection) was never actually exercised, since
nothing was ever tampered with in the first place. This report's
\emph{prediction} that the connection would fail to complete within 10
seconds was \textbf{incorrect}: because the attacker's diversion rule never
matched, the AH-wrapped packets simply passed through the attacker's
forwarding path untouched -- a transparent, non-tampering relay -- and the
connection completed normally. Rather than a hung connection, the real
result is arguably a stronger one for the defense: the legitimate connection
keeps working \emph{and} is immune to the attack, instead of merely being
denied service by IPsec's own failure mode. The 202 retransmits and reduced
throughput are a separate, secondary effect, not an attack or defense
artifact: this run's capture reported 5{,}587 packets dropped by
\texttt{tcpdump}'s own kernel buffer (out of $\sim$9{,}000 received),
indicating real CPU contention on the single physical machine running ARP
spoofing, software AH encryption/verification on every packet, NFQUEUE, and
packet capture simultaneously -- consistent with genuine, non-adversarial
packet loss under load rather than anything the attacker or the defense
caused deliberately. It is also worth noting explicitly that this run did
\emph{not} test AH's cryptographic rejection path (outcome a) at all, since
mechanism (b) pre-empted it; confirming outcome (a) specifically would
require diverting on IP protocol 51 instead of TCP flags and having the
rewriter attempt to blindly edit bytes inside the AH-protected region, which
is outside this report's scope.

\subsection{Defense C: Detection (Live WScale-Mismatch Monitor)}
\label{sec:defense-ids}

\paragraph{Process.} \texttt{defense/live\_ids.py} sniffs the bridge
interface directly (via \texttt{scapy}), tracking the WScale value seen on
the first SYN of every flow and comparing it against any later SYN sharing
the same sequence number, and separately flagging any data packet whose
window -- interpreted using the most recently observed scale for that flow
-- computes to under 16\,KiB. Unlike Defenses A and B, this is detect-only
and blocks nothing, so it was deliberately tested \emph{without} ARP
hardening or IPsec active, watching the same undefended deflate attack from
Section~\ref{sec:deflate} via
\texttt{sudo bash defense/run\_defense\_demo.sh ids}.

\paragraph{Expected Outcome.} The monitor should print one \texttt{[ALERT]}
line for the WScale mismatch on the tampered SYN, and further
\texttt{[ALERT]} lines for the anomalously small effective windows that
follow it during the throttled transfer -- both signatures already
confirmed present in the pcap analysis in Section~\ref{sec:deflate}.

\paragraph{Actual Outcome.} The detection logic had already been verified
offline (Process, above); it was then run live against the real, on-path
deflate attack from Section~\ref{sec:deflate}
(\texttt{sudo bash defense/run\_defense\_demo.sh ids}) and raised both alert
types predicted:
\begin{lstlisting}
[ALERT] 02:28:01  WScale mismatch on 10.0.0.1:40708 -> 10.0.0.2:9000
    (seq=1123647667): first saw 10, now 0 -- SYN tampered in flight

[ALERT] 02:28:01  Anomalous window on 10.0.0.1:40708 -> 10.0.0.2:9000:
    advertised=63 wscale=0 effective=63 bytes
[ALERT] 02:28:01  Anomalous window on 10.0.0.1:40708 -> 10.0.0.2:9000:
    advertised=63 wscale=0 effective=63 bytes
    ... (24,891 anomalous-window alerts total)
\end{lstlisting}
Exactly \textbf{1} mismatch alert fired, on the correct SYN, with the
correct before/after shift values (10 then 0) -- there is only one tampered
SYN in the whole run, so one alert is the correct count, not an
under-detection. The \textbf{24,891} anomalous-window alerts are consistent
with, not contradictory to, everything already established in
Section~\ref{sec:deflate}: a deflated connection is throttled to
$\sim$63-byte effective windows for its entire transfer, which (as computed
in Section~\ref{sec:deflate}) means several hundred thousand tiny segments
for a 20\,MB transfer -- \texttt{live\_ids.py}, unlike the pcap-based
detector used elsewhere in this report, has no packet cap, so it alerted on
every single one it saw during the run rather than a capped sample.

\paragraph{Does it work as expected?} \textbf{Yes, exactly.} Both signatures
the deck's Detection slide calls for -- a WScale mismatch between SYNs, and
anomalous window sizes -- fired correctly against real, live, on-the-wire
traffic during an actual attack, with the mismatch alert count matching the
true number of tampered SYNs exactly (one) and zero false positives of
either type.

% ---------------------------------------------------------------------------
\section{Assumptions}
\label{sec:assumptions}

The following assumptions underlie the attacks demonstrated in this report,
either because the attack's design depends on them or because the lab's
construction depends on them:

\begin{enumerate}
  \item \textbf{On-path position.} The attacker is already able to reach
        the same broadcast/Layer-2 segment as both the client and the
        server (no VLAN separation, no port isolation, no switch-level
        protections such as Dynamic ARP Inspection or 802.1X). This is the
        standard precondition for ARP spoofing and is assumed to already be
        satisfied, not obtained by some other exploit.
  \item \textbf{No transport/network-layer integrity protection.} Neither
        endpoint uses IPsec (AH/ESP) or any other mechanism that would
        authenticate the TCP header. TLS is explicitly \emph{not} a
        counter-measure here, since TLS protects only the application
        payload and never the TCP header the Window Scale option lives in.
  \item \textbf{Single-homed attacker relying on kernel forwarding.} The
        attacker has one network interface, not two; it relies on IP
        forwarding plus its own (unpoisoned) ARP resolution of the real peer
        to relay traffic back out the same interface, rather than bridging
        two separate physical links. This is the classic ``Ettercap-style''
        single-NIC MITM, and it is why the same logical packet is observed
        twice on the shared bridge (Section~\ref{sec:deflate}).
  \item \textbf{Window Scale is negotiated once and never re-verified.} The
        entire premise of Attacks 1 and 2 depends on the RFC~7323 behaviour
        that the shift value is fixed for the life of the connection after
        the handshake. If either endpoint's TCP stack ever re-advertised or
        re-validated the scale mid-connection, the desync would be
        self-correcting.
  \item \textbf{Both endpoints request window scaling.} If either side
        omits the \texttt{WScale} option from its SYN, both sides fall back
        to the unscaled 64\,KB window and there is no scale value left for
        the attacker to tamper with.
  \item \textbf{Endpoint TCP-stack behaviour is not fixed or fully
        predictable.} The original attack narrative implicitly treats the
        client's real window scale, and the visibility of an inflate
        attack's damage, as fixed quantities. In practice, Linux's
        receive-buffer autotuning and per-destination TCP metrics caching
        both meaningfully changed what was observed between otherwise
        identical runs (Section~\ref{sec:inflate}). Any claim of ``expected
        outcome'' for the inflate attack specifically should be understood
        as conditional on the receiving buffer actually being smaller than
        what the false window implies -- which is not guaranteed by the
        attack alone and had to be engineered explicitly in this lab via
        \texttt{SO\_RCVBUF}.
  \item \textbf{Lab timing is not representative of a real WAN.} The lab
        uses Linux network namespaces and a software bridge on one physical
        machine, giving near-zero RTT and very high effective bandwidth
        compared to a real network. This suppresses some effects (the
        inflate attack's overrun needed to be engineered around it) and
        exaggerates others (the deflate attack's packet-count explosion is
        partly a function of how fast one machine can loop packets to
        itself). The \emph{direction} of every observed effect should
        generalize to a real network; the specific magnitudes and timings
        should not be assumed to transfer directly.
  \item \textbf{Only the client's SYN is tampered, by design.} All three
        attacks target the client's plain SYN specifically
        (\texttt{--target client-syn}), matching the deck's own timing
        diagram, and leave the SYN-ACK untouched. The implementation also
        supports tampering the SYN-ACK or both packets
        (\texttt{--target server-synack}\,/\,\texttt{both}), but those
        configurations were not the ones demonstrated in this report.
  \item \textbf{The Slow Read fallback assumes a bounded server resource.}
        Attack 3 assumes the target server has some finite, exhaustible
        resource gated per-connection (here, a fixed-size worker thread
        pool). A server architecture without such a bound (e.g., a
        fully event-driven server with no per-connection worker limit)
        would not be denied service by this specific mechanism, though it
        may still suffer memory pressure from enough held-open connections.
  \item \textbf{Static ARP hardening assumes the real MAC addresses are
        already known and trustworthy at the time it is applied.}
        Defense~A (Section~\ref{sec:defense-arp}) pins each endpoint's real
        MAC into the other's neighbor table; it must be applied \emph{before}
        an attacker's first poison packet succeeds, and it requires some
        trusted way to learn the real MAC in the first place (here, reading
        it directly from each namespace's own interface, which a real switch
        would instead learn from legitimate DHCP/first-seen traffic under
        Dynamic ARP Inspection). It also does not scale automatically to
        hosts added after the fact without re-running the same process.
  \item \textbf{IPsec AH in this report uses a manually pre-shared key, not
        a production key-exchange protocol.} Defense~B
        (Section~\ref{sec:defense-ipsec}) demonstrates the authentication
        \emph{mechanism} (a shared key the attacker does not possess makes
        tampering detectable), not a deployable configuration; a real
        deployment would use IKE for key negotiation and rotation, which
        this report does not implement or evaluate.
  \item \textbf{The live IDS (Defense~C) is detection-only and assumes a
        human or automated system acts on its alerts.} It does not itself
        block, rate-limit, or otherwise interfere with the traffic it
        flags; it is a complement to Defenses~A and~B, not a substitute for
        either.
\end{enumerate}

% ---------------------------------------------------------------------------
\section{Summary}
\label{sec:summary}

\begin{table}[h]
\centering
\caption{Overall summary: expected vs.\ actual, all three attacks}
\begin{tabular}{p{2.7cm}p{4.3cm}p{4.3cm}p{2.6cm}}
\toprule
Attack & Expected & Actual & Matched? \\
\midrule
Deflate ($10\to0$) & Large throughput collapse, no packet loss & 99.9\% throughput drop, 0 retransmits, mismatch \& anomalous windows confirmed & Yes, exactly \\
Inflate ($\to14$), no pinned buffer & Retransmission storm / drops & No effect (throughput \emph{rose}, 0 retransmits) & No -- lab autotuning masked it \\
Inflate ($\to14$), pinned buffer & Retransmission storm / drops & 306--1398 retransmits across repeated runs, escalating & Yes, once corrected \\
Slow Read DoS & Pool exhaustion, legitimate requests denied & 20/50 connections held (exact match to pool size), \texttt{/health} denied during attack, recovered after & Yes, exactly \\
\bottomrule
\end{tabular}
\end{table}

Two of the three attacks (Deflate and Slow Read) matched their predicted
outcome on the first properly-executed attempt, with strong, unambiguous
evidence at both the socket level (\texttt{TCP\_INFO}) and the packet level
(pcap mismatch and anomalous-window detection). The Inflate attack required
one methodological correction -- explicitly pinning the client's real
receive buffer -- before its predicted failure mode (retransmission storm)
became observable, because the lab's fast, low-latency link let Linux's
normal receive-buffer autotuning hide the effect that a slower or more
resource-constrained real-world host would not have hidden. This divergence
and its resolution are documented in Section~\ref{sec:inflate} and are, in
themselves, a useful finding: an attacker relying solely on the ``inflate''
variant against a well-provisioned, high-bandwidth host may see no visible
effect at all, while the same attack against a resource-constrained host
(or over a slower/longer path) would be expected to succeed as originally
predicted.

\begin{table}[h]
\centering
\caption{Defense mechanisms: mechanism and observed status}
\begin{tabular}{p{2.6cm}p{5.2cm}p{4.4cm}p{1.9cm}}
\toprule
Defense & Mechanism & Observed effect on the attack & Status \\
\midrule
A: Static ARP (\S\ref{sec:defense-arp}) & Permanent neighbor entries reject forged ARP replies & Attacker never got on-path; 0 rewrites, 0 mismatches, full baseline throughput & \textbf{Confirmed working} \\
B: IPsec AH (\S\ref{sec:defense-ipsec}) & AH's protocol-51 wrapping stopped the diversion rule from matching the SYN at all & 0 rewrites, 0 mismatches, wscale untouched; transfer completed (not blocked as predicted -- see \S\ref{sec:defense-ipsec}) & \textbf{Confirmed working} \\
C: Live IDS (\S\ref{sec:defense-ids}) & Compares WScale across duplicate SYNs; flags sub-16\,KiB effective windows & 1 mismatch alert (exact match to the 1 tampered SYN), 24{,}891 anomalous-window alerts, 0 false positives & \textbf{Confirmed working} \\
\bottomrule
\end{tabular}
\end{table}

All three defenses are implemented as executable code
(\texttt{defense/harden\_arp.sh}, \texttt{defense/setup\_ipsec\_ah.sh},
\texttt{defense/live\_ids.py}, orchestrated by
\texttt{defense/run\_defense\_demo.sh}) and all three were run live against
the exact attack pipeline from Section~\ref{sec:deflate}, not merely
described. All three worked. Defenses A and B (static ARP, IPsec AH) both
defeated the attack completely -- zero rewrites, zero pcap mismatches in
either run -- though, notably, this report's own \emph{prediction} for
Defense B was partially wrong: it expected the connection to fail to
complete, when in fact it completed normally, simply immune to the
tampering (Section~\ref{sec:defense-ipsec} discusses why in detail, itself
a useful expected-vs-actual finding in the same spirit as
Section~\ref{sec:inflate}'s). Defense C (the live IDS) correctly raised
exactly one mismatch alert -- matching the true number of tampered SYNs
exactly -- plus tens of thousands of anomalous-window alerts consistent
with the deflate attack's already-established tiny-segment behaviour, with
no false positives, both offline against a synthetic test case and live
against real on-the-wire traffic.

Taken together with Sections~\ref{sec:deflate}--\ref{sec:slowread}, this
report now covers the full path the design report asked for: a
reproducible attack with measurable impact, at least one working defense
(in fact, three, covering blocking the MITM position, cryptographic
integrity, and detection), and an honest account of where this report's own
predictions -- for both the attacks and the defenses -- turned out to be
wrong, and why.

% ---------------------------------------------------------------------------
\section{References}
\label{sec:references}

The same six sources cited in the design report, annotated below with what
each was specifically used for in this implementation and report.

\begin{enumerate}[label={[\arabic*]}]
  \item RFC 7323, \emph{TCP Extensions for High Performance} (the Window
        Scale option). \url{https://www.rfc-editor.org/rfc/rfc7323} \\
        \textbf{Used for:} the entire attack's premise -- Kind~3/Length~3
        option format, the $2^{\text{shift}}$ multiplier, shift range
        $[0,14]$, and the specific behaviour exploited throughout this
        report: the option is sent only in the SYN and SYN-ACK and never
        re-verified. Directly informed \texttt{attacker/wscale\_attack.py}'s
        option parsing/rewrite logic (Section~\ref{sec:deflate}) and the
        effective-window calculation (\texttt{advertised\_window} $\times$
        $2^{\text{wscale}}$) used throughout
        \texttt{analysis/detect\_wscale\_mismatch.py} and
        \texttt{defense/live\_ids.py} (Section~\ref{sec:defense-ids}).
  \item RFC 9293, \emph{Transmission Control Protocol}.
        \url{https://www.rfc-editor.org/rfc/rfc9293} \\
        \textbf{Used for:} baseline TCP handshake, retransmission, and
        window/flow-control semantics against which every ``actual outcome''
        in this report was interpreted -- e.g., distinguishing the deflate
        attack's zero-retransmit self-throttling
        (Section~\ref{sec:deflate}) from the inflate attack's genuine
        retransmissions after real buffer overrun (Section~\ref{sec:inflate}).
  \item Fedora Magazine, \emph{TCP window scaling, timestamps and SACK}.
        \url{https://fedoramagazine.org/tcp-window-scaling-timestamps-and-sack/} \\
        \textbf{Used for:} a practical, implementation-level reference for
        how Linux specifically negotiates and applies window scaling,
        informing why the lab's default \texttt{snd\_wscale}/\texttt{rcv\_wscale}
        showed up as 10 rather than an RFC-example value, and why both
        endpoints falling back to unscaled 64\,KB windows (if either omits
        the option) matters as an assumption (Section~\ref{sec:assumptions},
        item 5).
  \item Qualys, \emph{Are you ready for slow reading?}
        \url{https://blog.qualys.com/vulnerabilities-threat-research/2012/01/05/slow-read} \\
        \textbf{Used for:} the design basis for the Slow Read fallback
        attack (Section~\ref{sec:slowread}) -- specifically the technique of
        advertising a tiny real receive window and reading the response a
        few bytes at a time to hold a server connection open indefinitely,
        which \texttt{slowread/slow\_read\_client.py} implements directly
        (small \texttt{SO\_RCVBUF}, paced \texttt{recv()} calls).
  \item \texttt{slowhttptest} tool and Slow Read wiki.
        \url{https://github.com/shekyan/slowhttptest} \\
        \textbf{Used for:} the reference implementation of the same attack
        class. \texttt{slowhttptest} itself was installed in the lab
        environment but was \emph{not} the tool actually invoked for the
        results in Section~\ref{sec:slowread}: a custom Python client
        (\texttt{slowread/slow\_read\_client.py}) was written instead, so
        that the pool-exhaustion target
        (\texttt{server/http\_server.py}'s bounded worker pool) and the
        attack client could be co-designed and instrumented together
        (e.g., the before/during/after \texttt{/health} checks that produced
        the evidence in Section~\ref{sec:slowread}). This is noted
        explicitly so the report does not imply \texttt{slowhttptest}'s own
        output was what was measured.
  \item CVE-2008-4609 (Sockstress / TCP persist-timer DoS), National
        Vulnerability Database.
        \url{https://nvd.nist.gov/vuln/detail/CVE-2008-4609} \\
        \textbf{Used for:} background context establishing Slow Read /
        zero-window abuse as a member of a broader, previously-disclosed
        class of low-bandwidth TCP state-exhaustion attacks, motivating why
        it was included as a fallback rather than treated as a novel
        technique in this report.
\end{enumerate}

\end{document}
